\documentclass{article}
\usepackage{etoolbox}
\usepackage{graphicx}

\ifdef{\myoutput}{}{\def\myoutput{concept}}
\ifdef{\myversion}{}{\def\myversion{1}}

\usepackage[output=\myoutput
,numberofversions=4
,version=\myversion
,seed=4
,randomizequestions=true
,randomizeanswers=true
,writeRfile=true
]{mcexam}

\usepackage[margin=0.95in]{geometry}
\usepackage{multicol}

\usepackage{fancyhdr,lastpage}
\pagestyle{fancy}
\fancyhf{}
\renewcommand{\headrulewidth}{1pt} 
\renewcommand{\footrulewidth}{1pt}
\lfoot{\mctheversion}
\rfoot{Page \thepage\ of \pageref{LastPage}}
\lhead{Quiz 5.1 \mctheversion \\ MAT 2000: Introduction to Statistics}
\rhead{Name Print): \hspace{2.5 in}}

\usepackage{calc}
\renewenvironment{setmcquestion}
{\begin{minipage}[t]{\linewidth-\labelwidth}}
	{\end{minipage}\par}
	
\renewenvironment{setmcquestioninstruction}
{\begin{minipage}{\textwidth}}
	{\end{minipage}}
	
\renewenvironment{setmcanswers}{}{}
\setlist[setmcquestions]{label=\mcquestionlabelfmt{*}.
	,ref=\mcquestionlabelfmt{*}
	,itemsep=0.5\baselineskip
	,topsep=1\baselineskip
}

\renewcommand\mcversionlabelfmt[1]{\arabic{#1}}

%copy-pasting from MS Word, use regex to find [\*]*[a-d]\. and replace with \answer 

\raggedbottom

\begin{document}
	Circle the correct answer for each question.  If you're unsure, you can write numbers 1, 2, 3, etc. next to the choices to get partial credit.\\[.5cm]
	
	\begin{mcquestions}
		\begin{mcquestioninstruction}
			The following table shows the responses from a sample of 680 people in the General Social Survey to the question, "Do you sometimes drink more than you think you should?"\\
			
			\begin{tabular}{llll}
				Gender & Yes                     & No  & Total  \\
				Male   & 151                     & 177 & 328    \\
				Female & 92                      & 260 & 352    \\
				Total  & 243                     & 437 & 680   
			\end{tabular}
		\end{mcquestioninstruction}
		\question What proportion of males responded that they drink more than they should?
		\begin{mcanswerslist}
			\answer 151/243
			\answer[correct] 151/328
			\answer 243/680
			\answer 151/680
		\end{mcanswerslist}
		
		\question[follow] Does there appear to be an association between gender and whether a person responds that they drink more than they should?
		\begin{mcanswerslist}
			\answer Yes, since a different number of people of each gender are in the sample.
			\answer Yes, since the percent that responded that they drink more than they should differs from 50%.
			\answer[correct] Yes, since the proportion who responded they drink more than they should differs between males and females.
			\answer No, there does not appear to be an association between these two variables.
		\end{mcanswerslist}
		
		\begin{mcquestioninstruction}
			A study published in 2007 by Christopher Johnson, professor of music education and music therapy at the University of Kansas, revealed that students in elementary schools who took private music lessons scored around 20 percent higher in math scores on standardized tests, compared to students who did not take private music lessons.  Use this to answer the following three questions:
		\end{mcquestioninstruction}
		
		\question What is the explanatory variable?
		\begin{mcanswerslist}
			\answer Music therapy
			\answer[correct] Whether a student took private music lessons
			\answer Math score on standardized test
			\answer Grade level
		\end{mcanswerslist}
		
		\question[follow] What is the response variable?
		\begin{mcanswerslist}
			\answer Music therapy
			\answer Whether a student took private music lessons
			\answer[correct] Math score on standardized test
			\answer Grade level
		\end{mcanswerslist}
		
		
		\question[follow] Which of the following could be a potential confounding variable?
		\begin{mcanswerslist}
			\answer Household income
			\answer Socio-economic status
			\answer Whether the parents hired a math tutor
			\answer[correct] All of these
		\end{mcanswerslist}
	\end{mcquestions}
\end{document}